\documentclass[11pt]{article}
\usepackage[T1]{fontenc}
\usepackage{lmodern}
\usepackage{amsmath,amssymb,amsthm,mathtools}
\usepackage[margin=1in]{geometry}
\usepackage{microtype}
\usepackage[colorlinks=true,linkcolor=blue,citecolor=blue,urlcolor=blue]{hyperref}
\hypersetup{pdftitle={Nonuniqueness for a reversible cubic reaction-diffusion system},pdfauthor={Kaiwen Jin}}
\newtheorem{theorem}{Theorem}[section]
\newtheorem{lemma}[theorem]{Lemma}
\theoremstyle{definition}
\newtheorem{definition}[theorem]{Definition}
\newcommand{\Om}{\Omega}
\newcommand{\R}{\mathbb R}
\newcommand{\dd}{\,dx}
\newcommand{\dt}{\,dt}
\newcommand{\ub}{\underline b}
\newcommand{\norm}[1]{\left\lVert#1\right\rVert}
\title{Nonuniqueness for a reversible cubic reaction--diffusion system}
\author{Kaiwen Jin\\\texttt{yui.ui.siki@gmail.com}}
\date{}
\begin{document}
\maketitle
\begin{abstract}
We construct two global renormalized solutions with the same nonnegative
finite-entropy initial data for the reversible cubic reaction
$A+2B\rightleftharpoons3B$ on a three-dimensional ball. The two diffusion
coefficients are equal. One species starts from zero, while the other has an
integrable $|x|^{-1}$ singularity. The zero-species solution coexists with a
second solution obtained from seeds whose initial masses tend to zero.
A compact self-similar subsolution keeps the second solution nonzero.
We prove its exact initial trace and the square-root Sobolev bounds needed
for the full renormalized formulation. Thus entropy dissipation and local
Lipschitz continuity of the reaction do not imply uniqueness for all
nonnegative finite-entropy initial data.
\end{abstract}

\section{Introduction and main result}

Reaction--diffusion systems describe concentrations transported by diffusion
and changed by local reactions. Their entropy inequalities often remain
available when the concentrations are too irregular to make the reaction
terms integrable. Renormalized solutions express the equations after smooth
cutoffs in concentration space. For entropy-dissipating reactions, Fischer
proved that these solutions coincide with a sufficiently regular comparison
solution whose concentrations are bounded above and bounded away from zero
uniformly \cite[Definition~2 and Theorem~7, preprint version]{Fischer2017}. The distinction between this
weak--strong statement and uniqueness of arbitrary renormalized solutions
is also discussed in \cite[Remark~1.4]{FFKT2023}.

Here we consider nonnegative initial concentrations of finite entropy,
allowing an unbounded concentration and an initially absent species. In this
class, a reversible cubic reaction already gives nonuniqueness. Let
$\Om=B_1(0)\subset\R^3$ and consider
\begin{equation}\label{eq:system}
 \begin{cases}
 a_t=\Delta a-b^2(a-b),\\
 b_t=\Delta b+b^2(a-b),
 \end{cases}
 \qquad \partial_\nu a=\partial_\nu b=0\quad\hbox{on }\partial\Om.
\end{equation}
The reaction $A+2B\rightleftharpoons3B$, with unit forward and reverse
rates, has reaction vector
\begin{equation}\label{eq:reaction}
 f(a,b)=\bigl(-b^2(a-b),\ b^2(a-b)\bigr).
\end{equation}
It preserves the nonnegative quadrant and satisfies, for every $a,b>0$,
\begin{equation}\label{eq:entropy-sign}
 f_1(a,b)\log a+f_2(a,b)\log b
   =-b^2(a-b)(\log a-\log b)\le0.
\end{equation}
For entropy variables $\log u_i+\mu_i$, where the $\mu_i$ are fixed real
constants, this gives the entropy sign with $\mu_1=\mu_2=0$.

We use the following precise solution class. As usual, $0\log0=0$.
\begin{definition}\label{def:renormalized}
A nonnegative vector $u=(u_1,u_2)=(a,b)$ is a global renormalized solution
of \eqref{eq:system} with initial data $u_0$ if, for every $T<\infty$,
\begin{equation}\label{eq:class}
 u_i\in L^\infty(0,T;L^1(\Om)),\qquad
 \sqrt{u_i}\in L^2(0,T;H^1(\Om)),\quad i=1,2,
\end{equation}
and the following identity holds for every smooth
$\xi:[0,\infty)^2\to\R$ with compactly supported derivative $D\xi$, every
smooth $\psi$ on $\overline\Om\times[0,T]$, and almost every terminal time $T$:
\begin{align}\label{eq:renormalized}
 &\int_\Om\xi(u(T))\psi(T)\dd-\int_\Om\xi(u_0)\psi(0)\dd
       -\int_0^T\!\int_\Om\xi(u)\partial_t\psi\dd\dt\notag\\
 &\quad=-\sum_{i,j=1}^2\int_0^T\!\int_\Om
             \psi\xi_{ij}(u)\nabla u_i\cdot\nabla u_j\dd\dt
       -\sum_{i=1}^2\int_0^T\!\int_\Om
             \xi_i(u)\nabla u_i\cdot\nabla\psi\dd\dt\notag\\
 &\qquad\quad+\sum_{i=1}^2\int_0^T\!\int_\Om
             \xi_i(u)f_i(u)\psi\dd\dt.
\end{align}
Here $\nabla u_i=2\sqrt{u_i}\nabla\sqrt{u_i}$ in the gradient products.
\end{definition}

\begin{theorem}\label{thm:main}
Set $C=10000$ and
\begin{equation}\label{eq:data}
 a_0(x)=\frac C{|x|}\quad(x\ne0),\qquad b_0(x)=0,
\end{equation}
with any finite nonnegative value assigned to $a_0(0)$. These data satisfy
\[
 \sum_{i=1}^2\int_\Om u_{i,0}(1+|\log u_{i,0}|)\dd<\infty.
\]
System \eqref{eq:system} has at least two different global renormalized
solutions with these initial data. If $m$ is the Neumann heat evolution of
$a_0$, one solution is $(m,0)$. Another is $(m-b,b)$, where
\begin{equation}\label{eq:nonzero-bound}
 0\le b\le\frac m2\quad(t>0),\qquad
 b(t,x)\ge t^{-1/2}\left(1-\frac{|x|^2}{t}\right)_+^2
       \quad\left(0<t\le\frac14\right).
\end{equation}
Both solutions have the strong $L^1$ initial trace $(a_0,0)$ and are regular
on every compact positive-time interval.
\end{theorem}

The polynomial vector field is locally Lipschitz, and its boundary signs
are $f_1(0,b)=b^3\ge0$ and $f_2(a,0)=0$. Thus the theorem gives an instance
of an entropy-dissipating system with positive constant diffusivities in
the usual finite-entropy comparison class. The singular initial
concentration in \eqref{eq:data} is part of the construction.
Theorem~\ref{thm:main} gives a negative answer to Problem~482 in the AIM
collection~\cite{AIM2026}, for its stated class of nonnegative finite-entropy
data.

Equal diffusion is useful here: the total concentration $m=a+b$ solves a
linear heat equation. It therefore bounds the nonlinear evolution globally,
while its initial singularity permits a nonzero catalyst concentration to
grow from a sequence of seeds with vanishing mass. The main point at the
initial time is to show that this limiting concentration still has trace
zero. A localized estimate away from the singular point proves that trace;
the entropy estimate then gives \eqref{eq:class}.

\section{A heat profile and a compact subsolution}

Let $S(t)$ denote the Neumann heat semigroup on $\Om$, and put
$m(t)=S(t)a_0$. Since $a_0\in L^2(\Om)$, positivity, strong continuity and
positive-time regularity of the heat semigroup apply
\cite[Example~5.2]{SchnaubeltLinear}. In particular $m$ is bounded on each
$[\delta,T]\times\overline\Om$, $\delta>0$, and
\begin{equation}\label{eq:heat-trace}
 m(t)\longrightarrow a_0\text{ in }L^1(\Om),\qquad
 \int_\Om m(t)\dd=2\pi C.
\end{equation}
The following explicit bounds control the behavior near the initial time.

\begin{lemma}\label{lem:heat}
For $t>0$ and $x\ne0$,
\begin{equation}\label{eq:heat-bounds}
 \frac C{\sqrt{6t+|x|^2}}\le m(t,x)
       \le\frac C{|x|}+\frac C2|x|^2+3Ct.
\end{equation}
Consequently, for each finite $T$,
\begin{equation}\label{eq:small-ball}
 \lim_{\rho\downarrow0}\sup_{0<t\le T}
          \int_{B_{2\rho}}m(t,x)\dd=0.
\end{equation}
\end{lemma}
\begin{proof}
For $\eta>0$ set $m_{0,\eta}=C(|x|^2+\eta)^{-1/2}$ and
$m_\eta=S(t)m_{0,\eta}$. The function
$\ell_\eta=C(6t+|x|^2+\eta)^{-1/2}$ satisfies
\[
 (\partial_t-\Delta)\ell_\eta
       =-3C|x|^2(6t+|x|^2+\eta)^{-5/2}\le0,
 \qquad\partial_\nu\ell_\eta\le0.
\]
Its initial value is $m_{0,\eta}$, so $\ell_\eta\le m_\eta$ by comparison.
For the upper bound, define
\[
 w_\eta=C(|x|^2+\eta)^{-1/2}
             +\frac{C|x|^2}{2(1+\eta)^{3/2}}.
\]
Then $\partial_\nu w_\eta=0$ on $\partial\Om$, $w_\eta\ge m_{0,\eta}$,
and
\[
 \Delta w_\eta=-3C\eta(|x|^2+\eta)^{-5/2}
                  +3C(1+\eta)^{-3/2}\le3C.
\]
Hence $m_\eta\le w_\eta+3Ct$. These comparisons follow by testing the
positive part of the heat-equation difference; the boundary term for
$\ell_\eta-m_\eta$ is nonpositive. They only use strong $L^2$ initial
traces of the regularized data. Letting $\eta\downarrow0$ proves
\eqref{eq:heat-bounds}, since $m_{0,\eta}\to a_0$ in $L^2$. Its upper
bound has integral $O(\rho^2)+O_T(\rho^3)$ on $B_{2\rho}$, which proves
\eqref{eq:small-ball}.
\end{proof}

The lower bound in \eqref{eq:heat-bounds} has the same parabolic scale as
the following compactly supported function.
\begin{lemma}\label{lem:subsolution}
For $0<t\le T_0:=1/4$, let
\[
 \ub(t,x)=t^{-1/2}\Phi(x/\sqrt t),\qquad
 \Phi(y)=(1-|y|^2)_+^2.
\]
Then $0\le\ub\le m/2$, its support is contained in $\overline{B_{1/2}}$,
and it is a
weak subsolution of
\begin{equation}\label{eq:scalar}
 b_t=\Delta b+b^2(m-2b)
\end{equation}
with zero Neumann boundary flux.
\end{lemma}
\begin{proof}
Write $s=|y|^2$ and
$\mathcal L=\Delta_y+\frac12 y\cdot\nabla_y+\frac12$. Direct calculation
on $s<1$ gives
\[
 \mathcal L\Phi=-\frac{23}{2}+17s+\frac52s^2,
 \qquad(\partial_t-\Delta)\ub=-t^{-3/2}\mathcal L\Phi.
\]
Also $m(t,x)\ge t^{-1/2}C/\sqrt{6+s}$ on the support of $\ub$. Thus the
desired subsolution inequality follows from
\begin{equation}\label{eq:profile-ineq}
 -\mathcal L\Phi\le\Phi^2
                \left(\frac C{\sqrt{6+s}}-2\Phi\right).
\end{equation}
For $s\le3/4$, the left side is at most $23/2$ and the right side is at
least $(C/3-2)/256>23/2$. For $3/4\le s<1$, one has
$\mathcal L\Phi\ge85/32>0$, while the right side is nonnegative.
Moreover $2\Phi\le2<C/\sqrt7$, proving $\ub\le m/2$.
At $s=1$ both $\Phi$ and its first derivatives vanish. The extension by
zero therefore has no surface measure in its distributional Laplacian,
and \eqref{eq:profile-ineq} proves the weak inequality across the moving
front. The support radius is $\sqrt t\le1/2$, so the outer boundary flux
is zero.
\end{proof}

\section{A second solution and its initial trace}

For $0<\varepsilon<T_0$, start \eqref{eq:system} at time $\varepsilon$ with
\begin{equation}\label{eq:seed-data}
 (a_\varepsilon,b_\varepsilon)(\varepsilon)
       =\bigl(m(\varepsilon)-\ub(\varepsilon),\ub(\varepsilon)\bigr).
\end{equation}
These data belong to $W^{2,p}(\Om)^2$ with zero Neumann trace for every
finite $p$. Indeed, the compact seed is $C^1$, has bounded piecewise second
derivatives, and vanishes near $\partial\Om$. Bounded-data semilinear
Neumann theory gives a maximal regular solution
\cite[Example~3.7 and Theorem~3.9]{SchnaubeltNonlinear}.
Its total is $m$, by uniqueness of the scalar heat equation. Thus
$b_\varepsilon$ solves \eqref{eq:scalar}.

On each interval bounded away from zero the scalar reaction is locally
Lipschitz in $b$. Both $0$ and $m/2$ solve \eqref{eq:scalar}; comparison
and Lemma~\ref{lem:subsolution} yield
\begin{equation}\label{eq:approx-bounds}
 0\le b_\varepsilon(t)\le m(t)/2\quad(t\ge\varepsilon),\qquad
 b_\varepsilon(t)\ge\ub(t)\quad(\varepsilon\le t\le T_0).
\end{equation}
For the subsolution comparison, testing the positive part of the difference
gives its squared $L^2$ norm a differential inequality with a finite local
Lipschitz constant. No bound on this constant as $\varepsilon\downarrow0$
is needed. Since $a_\varepsilon=m-b_\varepsilon$, the bound by $m$ excludes
finite-time blowup of either component. Each seeded solution is global.

We spell out the positive-time compactness used to remove the seed time.
Fix $0<\delta<T$ and $\varepsilon<\delta/2$. The states and their polynomial
reaction terms are uniformly bounded on $[\delta/2,T]$. Put
$X=L^6(\Om)^2$, let $A=\operatorname{diag}(\Delta_N,\Delta_N)$, and write
$X_\theta=D_A(\theta,6)$ for its real interpolation spaces. The mild formula
and analytic smoothing imply
\begin{align*}
 \norm{(a_\varepsilon,b_\varepsilon)(\delta)}_{X_{2/3}}
 &\le K_\delta\norm{(a_\varepsilon,b_\varepsilon)(\delta/2)}_X\\
 &\quad+K_\delta\int_{\delta/2}^\delta(\delta-s)^{-2/3}
                 \norm{f(a_\varepsilon(s),b_\varepsilon(s))}_X\,ds.
\end{align*}
The right side is uniform in $\varepsilon$. The embedding
$X_{2/3}\hookrightarrow X_{1/2}$ is compact and
$X_{1/2}\hookrightarrow C(\overline\Om)^2$ is continuous; see
\cite[Proposition~2.13, Theorem~2.14 and Example~3.7]{SchnaubeltNonlinear}.
Extract a convergent subsequence at $\delta$, then use continuous dependence
for the bounded positive-time problem. Its limit remains bounded by $m$,
so it extends to all later times. Nested subsequences at times decreasing
to zero agree on overlaps and give one global regular solution
\begin{equation}\label{eq:limit}
 (a,b)=(m-b,b),\qquad0\le b\le m/2,
 \qquad b\ge\ub\quad(0<t\le T_0).
\end{equation}
The convergence at each positive time is uniform on $\overline\Om$.

\begin{lemma}\label{lem:trace}
The solution in \eqref{eq:limit} satisfies
$b(t)\to0$ and $a(t)\to a_0$ in $L^1(\Om)$ as $t\downarrow0$.
\end{lemma}
\begin{proof}
Fix $0<\rho<1/4$ and a smooth $0\le\chi_\rho\le1$ which is zero on
$B_\rho$ and one outside $B_{2\rho}$. It is constant near $\partial\Om$.
For each finite $T$, Lemma~\ref{lem:heat} bounds $m$ by some
$M_{\rho,T}$ on $0<t\le T$, $|x|\ge\rho$. Set
$Y_\varepsilon(t)=\int_\Om b_\varepsilon(t)\chi_\rho\dd$.
Since $b_\varepsilon^2(m-2b_\varepsilon)\le m^2b_\varepsilon$,
\[
 Y_\varepsilon'(t)\le K_\rho+M_{\rho,T}^2Y_\varepsilon(t),\qquad
 K_\rho=\tfrac12\norm{\Delta\chi_\rho}_\infty\int_\Om a_0\dd.
\]
The diffusion term is $\int b_\varepsilon\Delta\chi_\rho$; both boundary
terms vanish. If $\varepsilon<\rho^2$, the seed is supported in $B_\rho$,
so $Y_\varepsilon(\varepsilon)=0$. Hence
$Y_\varepsilon(t)\le K_\rho t\exp(M_{\rho,T}^2t)$.
Pass to the limit at fixed $t>0$ to obtain
\[
 \int_\Om b(t)\dd\le\frac12\int_{B_{2\rho}}m(t)\dd
                    +K_\rho t\exp(M_{\rho,T}^2t).
\]
First let $t\downarrow0$ with $\rho$ fixed, and then let $\rho\downarrow0$
using \eqref{eq:small-ball}. This proves the zero trace of $b$. The trace
of $a=m-b$ follows from \eqref{eq:heat-trace}.
\end{proof}

\section{Entropy and the renormalized identity}

Let $h(s)=s\log s-s+1$ for $s>0$, with $h(0)=1$, and put
$H(t)=\int_\Om[h(a(t))+h(b(t))]\dd$.
Since $a+b=m$, one has
\begin{equation}\label{eq:entropy-upper}
 0\le h(a)+h(b)\le h(m)+1.
\end{equation}
For any fixed finite $T$, \eqref{eq:heat-bounds} dominates the right side
on $0<t\le T$ by a constant multiple of
$1+|x|^{-1}\log(2+|x|^{-1})$, an integrable function on the ball.
Lemma~\ref{lem:trace}, continuity of $h$ and dominated convergence therefore
give
\begin{equation}\label{eq:entropy-trace}
 H(t)\longrightarrow H_0:=\int_\Om[h(a_0)+h(0)]\dd<\infty.
\end{equation}
For clarity, apply dominated convergence along an almost-everywhere
subsequence of any sequence converging to zero in time. Every such sequence
has this subsequence by the strong $L^1$ trace, so the full limit follows.

For $\eta>0$ use the smooth entropy $h_\eta(s)=h(s+\eta)$ on
$[\delta,T]$, where $\delta>0$. Write $H_\eta(t)=\sum_i\int h(u_i(t)+\eta)$.
The equations and the Neumann condition give
\begin{align}\label{eq:regularized-entropy}
 H_\eta(T)&+\sum_{i=1}^2\int_\delta^T\!\int_\Om
                      \frac{|\nabla u_i|^2}{u_i+\eta}\dd\dt\notag\\
 &+\int_\delta^T\!\int_\Om b^2(a-b)
              \bigl[\log(a+\eta)-\log(b+\eta)\bigr]\dd\dt
       =H_\eta(\delta).
\end{align}
The last integrand is nonnegative because $a\ge b$. This calculation is
valid even where $b=0$. It bounds $\sqrt{u_i+\eta}$ in
$L^2(\delta,T;H^1(\Om))$, while
$|\sqrt{u_i+\eta}-\sqrt{u_i}|\le\sqrt\eta$ gives strong $L^2$ convergence.
Weak compactness and lower semicontinuity, followed by $\eta\downarrow0$,
yield
\[
 4\sum_{i=1}^2\int_\delta^T\!\int_\Om|\nabla\sqrt{u_i}|^2\dd\dt
       \le H(\delta)-H(T).
\]
Now let $\delta\downarrow0$ and use \eqref{eq:entropy-trace}. Together
with $a+b=m$ and its conserved mass, this proves \eqref{eq:class} for every
finite $T$.
The gradient identity used in Definition~\ref{def:renormalized} follows as
well: test $\nabla u_i=2\sqrt{u_i+\eta}\nabla\sqrt{u_i+\eta}$ on a
fixed positive-time strip, and pass to the limit using strong $L^2$
convergence of the first factor and weak $L^2$ convergence of the second.


\begin{proof}[Proof of Theorem~\ref{thm:main}]
The initial entropy in \eqref{eq:data} is finite, since its radial integral
is a constant multiple of $\int_0^1 r[1+\log(C/r)]\,dr$.
The construction above gives a global solution $(m-b,b)$ with the required
initial trace and square-root Sobolev bounds. It remains to check the full
identity \eqref{eq:renormalized} at the initial time.

For fixed $\delta>0$, the smooth chain rule gives that identity on
$[\delta,T]$. On the common compact support of $D\xi$ and $D^2\xi$, the
components $u_i$ are bounded by some $R$. The identity
$\nabla u_i=2\sqrt{u_i}\nabla\sqrt{u_i}$ then bounds every Hessian term by
a constant times
$|\nabla\sqrt{u_i}|\,|\nabla\sqrt{u_j}|$ and every flux term by a constant
times $|\nabla\sqrt{u_i}|$. These terms are integrable on $(0,T)\times\Om$.
The cutoff reaction term is bounded. Also $\xi$ is globally Lipschitz, so
the strong $L^1$ initial trace gives $\xi(u(\delta))\to\xi(u_0)$ in $L^1$.
Letting $\delta\downarrow0$ therefore proves \eqref{eq:renormalized} for
every permitted $\xi$ and $\psi$, in fact for every $T>0$.

The second solution $(m,0)$ has zero reaction, the same strong initial
trace, and the same entropy dominator \eqref{eq:entropy-upper}. The
regularized entropy calculation \eqref{eq:regularized-entropy} applies to
it with zero reaction term, giving its Fisher bound without testing
$\log0$. The identical chain-rule passage proves
\eqref{eq:renormalized} for this solution as well.
Finally $b\ge\ub>0$ on $0<t<T_0$, $|x|<\sqrt t$, a set of positive
space-time measure. The two renormalized solutions are different.
\end{proof}

\begin{thebibliography}{9}
\bibitem{Fischer2017}
J.~Fischer, \emph{Weak--strong uniqueness of solutions to entropy-dissipating
reaction--diffusion equations}, Nonlinear Analysis \textbf{159} (2017),
181--207. \href{https://doi.org/10.1016/j.na.2017.03.001}{doi:10.1016/j.na.2017.03.001}.
The theorem numbering used here is from the preprint version
\href{https://arxiv.org/abs/1703.00730v1}{arXiv:1703.00730v1}.

\bibitem{FFKT2023}
K.~Fellner, J.~Fischer, M.~Kniely and B.~Q.~Tang,
\emph{Global renormalised solutions and equilibration of reaction--diffusion
systems with nonlinear diffusion}, Journal of Nonlinear Science
\textbf{33} (2023), article 66.
\href{https://doi.org/10.1007/s00332-023-09926-w}{doi:10.1007/s00332-023-09926-w}.

\bibitem{SchnaubeltLinear}
R.~Schnaubelt, \emph{Evolution Equations}, lecture notes, Karlsruhe Institute
of Technology, version March 19, 2026.
\url{https://iana.math.kit.edu/downloads/iana3/schnaubelt/Skripten/evgl-skript.pdf}.

\bibitem{SchnaubeltNonlinear}
R.~Schnaubelt, \emph{Nonlinear Evolution Equations}, lecture notes,
Karlsruhe Institute of Technology, version October 8, 2025.
\url{https://iana.math.kit.edu/downloads/iana3/schnaubelt/Skripten/nevgl-skript.pdf}.
\bibitem{AIM2026}
S.~Brunton, M.~J.~Colbrook, M.~V.~de~Hoop, G.~Stepaniants, A.~Townsend and
R.~Ward, \emph{AIM---Open Applied Problems}, GitHub repository (2026).
\href{https://github.com/MColbrook/AIM/blob/a153f469ee0ea64e876b2ad97c79d252c19a3e22/problems/482-entropy-reaction-diffusion-renormalized-uniqueness.md}{Problem~482},
revision \texttt{a153f469ee0ea64e876b2ad97c79d252c19a3e22}.
\end{thebibliography}

\section*{AI assistance}
OpenAI Codex assisted with developing the construction, checking calculations,
locating references and drafting this manuscript. Independent mathematical
review is requested.
\end{document}
